\section{Implementation Details} \label{sec:implementation}

We have implemented the proposed algorithms within the LAGraph project~\cite{lagraph}, a collection of linear algebra-based graph processing algorithms built on top of \ssgrb \cite{10.1145/3577195}. This section describes the key implementation aspects.

% The actual implementation involves an additional semiring due to the asymmetric nature of the definition. The original semiring $\paths = \langle \rpaths, \rpaths, \mathbb{B} = \{ 0, 1 \}, \cup, \textbf{fst}, \emptyset \rangle$ uses $\textbf{fst} : \rpaths \times \mathbb{B} \rightarrow \rpaths$.
%
% However, it is convenient to introduce another semiring $\overline{\paths} = \langle \rpaths, \mathbb{B} = \{ 0, 1 \}, \rpaths, \cup, \textbf{snd}, \emptyset \rangle$, which uses $\textbf{snd} : \mathbb{B} \times \rpaths \rightarrow \rpaths$, defined as ??. 
%
% If we use these two semirings and the corresponding multiplicative operations $\otimes$ and $\overline{\otimes}$, we obtain:
%
% \begin{equation*}
%   (A \otimes B)^T = (B^T \overline{\otimes} A^T)
% \end{equation*}
%
% for matrices $A$ and $B$ over $D_{out}$.
%
% Using this observation, we can avoid the expensive product transposition on line ?? in Algorithm \ref{algo:semiring-rpq} and use the cheaper $(N^a)^T$ operation instead:
%
% \begin{equation*}
%   ((M^T \otimes N^a)^T \otimes G^a) \leadsto ((N^a)^T \overline{\otimes} M \otimes G^a).
% \end{equation*}
% \rs{should this dot be there?}
%
% In addition, \ssgrb imposes restrictions on user-defined semirings. Most importantly, user-defined semirings must have fixed-size values to work with built-in automatic allocators.
% Therefore, the implementation of $\paths$ in \ssgrb relies on external memory management when paths become too long or when too many paths appear in a single $M_{qv}$ $\rpaths$ element.
% We introduce two constants, $L_{fast}$ and $C_{fast}$, to control the maximum path length and the maximum cardinality of a $\rpaths$ element without external memory management.
%
% If a path $\pi \in M_{qv}$ with $|\pi| > L_{fast}$ appears, a new linked list representing this long path is allocated in a supplementary memory arena.
% Similarly, if $|M_{qv}| > C_{fast}$ for some $q, v$, an extensible linked list is allocated.
%
% These two values provide a trade-off between peak memory consumption and execution speed.
% In practice, these constants should be adjusted for each dataset individually.
%
% Finally, we ship the semiring-based algorithm implementation as a standalone executable file.
% A preprocessed graph, decomposed into a set of Boolean Matrix Market matrices, and a set of regular path query patterns described in a SPARQL-like query language are consumed as input and executed one by one.
% The CLI flag ?? \texttt{--semantics} can be used to switch between all-simple-path and all-trail semantics.
% At the same time, \texttt{--fast-path-length} controls $L_{fast}$, and \texttt{--fast-path-cardinality} defines the $C_{fast}$ constant.

The main engineering limitation is that \ssgrb user-defined semiring values must have a fixed size to work with built-in automatic allocators.
At the same time, an element of the path semiring may contain paths of different lengths and may accumulate an unbounded number of paths during evaluation.
To fit these variable-size objects into the fixed-size semiring interface, the implementation uses a compact in-place representation for common cases and falls back to explicit memory management only when this representation overflows.

We introduce two constants, $L_{fast}$ and $C_{fast}$, to control the maximum path length and the maximum cardinality of a $\rpaths$ element that can be stored without external memory management.
As long as all paths in an element are no longer than $L_{fast}$ and the element contains at most $C_{fast}$ paths, the value is stored directly inside the fixed-size semiring object.
This fast representation avoids extra allocations and is intended to cover the typical case.

If a path $\pi \in M_{qv}$ with $|\pi| > L_{fast}$ appears, a new linked list representing this long path is allocated in a supplementary memory arena.
Similarly, if $|M_{qv}| > C_{fast}$ for some $q, v$, an extensible linked list is allocated.
These fallback structures keep the semiring value itself fixed-size: it stores only the inline data or references to the arena-managed objects.
We used $L_{fast} = 8$ and $C_{fast} = 1$, empirically chosen to achieve the best performance in evaluation.

These two values provide a trade-off between peak memory consumption and execution speed.
In practice, these constants should be adjusted for each dataset individually.

Finally, we ship the semiring-based algorithm implementation as a standalone executable file.
A preprocessed graph, decomposed into a set of Boolean Matrix Market matrices, and a set of regular path query patterns described in a SPARQL-like query language are consumed as input and executed one by one.
The CLI flag \texttt{--semantics} can be used to switch between all-simple-path, all-shortest-path and all-trail semantics.
